\documentclass[10pt]{article}

\usepackage[utf8]{inputenc}

\usepackage[T1]{fontenc}

\usepackage{amsmath}

\usepackage{amsfonts}

\usepackage{amssymb}

\usepackage[version=4]{mhchem}

\usepackage{stmaryrd}

\usepackage{mathrsfs}

\usepackage{graphicx}

\usepackage[export]{adjustbox}

\graphicspath{ {./images/} }

\usepackage{bbold}



\begin{document}

ЭКСТРЕМУМ - значение непрерывной функции, являющееся максимумом или минимумом (см. Максияужа и минимума точки). Термин «Э.» употребляется также при изучении наибольших и наименьших значений функционалов в вариационном исчислении. А. Б. Иванов.



ЭКСЦЕНТРИСИТЕТ - число, равное отношению расстояния от любой точки конического сечения до данной точки (фокуса) к расстоянию от той же точки до данной прямой (директрисы). Два конич. сечения, имеющие равные Ә., подобны между собой. Для әллипса Ә. $e<1$ (для окружности $e=0$ ), для гиперболы $e>1$, для параболы $e=1$. Для эллипса и гиперболы Э. можно определить как отношение расстояний между фокусами к длине бо́лышей оси. $А$. Б. Иванов. ЭКСЦЕССА КОЭФФИЦИЕНТ, ә к с ц е с с, - скалярная характеристика островершинности графика плотности вероятности унимодального распределения, к-рую используют в качестве нек-рой меры отклонения рассматриваемого распределения от нормального. Э. к. $\gamma_{2}$ определяется по формуле



\[

\gamma_{2}=\beta_{2}-3

\]



где $\beta_{2}=\mu_{4} / \mu_{2}^{2}$ есть 2 -й коэффициент Пирсона, $\mu_{2}^{-}$и $\mu_{4}-$ 2 -й и 4-й центральные моменты вероятностного распределения. В терминах семиинвариантов (кумулянтов) $x_{2}$ и $x_{4}$ второго п четвертого порядков Э. к. имеет вид



\[

\gamma_{2}=\frac{x_{4}}{x_{2}^{2}}

\]



Если $\gamma_{2}=0$, то говорят, что плотность вероятности распределения имеет н о р м а л ь н й ибо для нормального распределения Ә. к. $\gamma_{2}=0$. В случае, если Э. к. $\gamma_{2}>0$, то говорят, что плотность вероятности имеет по ло жи те ль ны й ә к с ц е с с, что, как правило, соответствует тому, что график плотности рассматриваемого распределения в окрестности моды имеет более острую и более высокую вершину, чем нормальная кривая. В случае, когда $\gamma_{2}<0$, говорят об о т р и ца т е ль н о м ә к с д е с с е плотности, при әтом плотность вероятности имеет в окрестности моды более низкую и плоскую вершину, чем плотность нормального закона.



Если $X_{1}, X_{2}, \ldots, X_{n}$ - независимые случайные величины, подчиняющиеся одному и тому же непрерывному вероятностному закону, то статистика



\[

\hat{\gamma}_{2}=\frac{1}{n\left(s^{2}\right)^{2}} \sum_{i=1}^{n}\left(X_{i}-X_{\theta}\right)^{4}-3

\]



наз. вы бо ро чным к о бффициентом әксц е с с а, где



\[

X=\frac{1}{n} \sum_{i=1}^{n} X_{i}, \quad s^{2}=\frac{1}{n} \sum_{i=1}^{n}\left(X_{i}-X_{0}\right)^{2}

\]



Выборочный Э. к. $\hat{\gamma}_{2}$ используется в качестве точечноӥ статистич, оценки Э. к. $\gamma_{2}$, когда закон распределения случайной величины $X_{i}$ неизвестен. В случае нормальной распределенности случайных величин $X_{1}, \ldots, X_{n}$ выборочный Ә. к. $\hat{\gamma}_{2}$ асимптотически нормально распределен при $n \rightarrow \infty$ с параметрами



\[

E \hat{\gamma}_{2}=-\frac{6}{n+1}

\]



и



\[

\mathrm{D} \hat{\gamma}_{2}=\frac{24 n(n-2)(n-3)}{(n+1)^{2}(n+3)(n+5)}=\frac{24}{n}\left[1-\frac{225}{15 n+24}+O\left(\frac{1}{n^{3}}\right)\right] .

\]



Именно поэтому, если наблюденное значение выборочного Э. к. $\hat{\gamma}_{2}$ существенно отличается от 0 , то следует признать, что распределение случайной величины $X_{i}$ отлично от нормального, чем и пользуются на практике прп проверке гипотезы $H_{0}: \gamma_{2} \neq 0$, пршнятие к-рой равносильно признанию отклонения распределения случайной величины $X_{i}$ от нормального. Лum.: [1] К е н д а л л . М. Д ж., Сть ю а р т А., Теория распределений, пер. с англ., М., 1966; [2] К р а м е р Г., Математические методы статистики, пер. с англ., 2 изд., М., 1975; тематической статистики, 3 изд., м., 1983. М. М. Никулик.



ӘКСЦЕССИВНАЯ ФУНКЦИЯ Для марковского процесса - аналог неотрицательной супөргармонической функции.



Пусть в измеримом пространстве $(E, \mathscr{B})$ задана однородная марковская цепь с вероятностями перехода за один шаг $P(x, B)(x \in E, B \in \mathscr{B})$. Измеримая относи-



\includegraphics[max width=\textwidth, center]{2023_07_09_ec7bb519de6655bea482g-1}

н $о$ й ф у н $к$ ц и $ө$ й для этой цепи, если



\[

\int_{E} f(y) \boldsymbol{P}(x, d y) \leqslant f(x)

\]



всюду в $E$. Для цепи с не более чем счетным множеством состояний среди Ә. ф. есть отличные от констант тогда И только тогда, когда хоть одно из состояний невозвратно.



При заданном в $(E, \mathscr{B})$ однородном марковском процессе $X=\left(x_{t}, \zeta, \mathscr{F}_{t}, P_{x}\right)$ с переходной функцией $P(t, x$, $B)$ определение Э. ф. несколько усложнится. Множество $B$ относят к $\sigma$-алгебре $\overline{\mathscr{B}}$, если для любой конечной меры $\mu$ на $\mathscr{B}$ найдутся такие множества $B_{\mu}^{1}$ и $B_{\mu}^{2}$, что $B_{\mu}^{1} \subset B \subset B_{\mu}^{2}$ и $\mu\left(B_{\mu}^{2} \mid B_{\mu}^{1}\right)=0$. Функция $f_{:} \quad E \longrightarrow$ $\longrightarrow[0, \infty]$ наз. ә к с це с с и в о ой фу укци е й, если она $\mathscr{B}$-измерима и при $t \geqslant 0$ удовлетворяет всюду в $E$ соотношениям



\[

\boldsymbol{P t}_{f}(x) \equiv \int f(y) \boldsymbol{P}(t, x, d y) \leqslant f(x)

\]



и



\[

f(x)=\lim _{s_{\swarrow} \swarrow 0} P^{s} f(x) .

\]



Для части винеровского процесса в нек-рой области $E \subset \mathbb{R}^{n}$ (см. Функционал от марковского процесса) класс Э. Ф. совпадает с классом неотрицательных супергармонических в $E$ функций, пополненным функцией $f(x) \equiv \infty$.



В случае стандартного процесса $X$ в локально компактном сепарабельном пространстве $E$, для Э. ф. $f(x)$ всюду в $E$ выполняется неравенство



\[

M_{x} f\left(x_{\tau}\right) \leqslant f(x)

\]



где $\tau$-марковский момент, $M_{x}$-математическое ожидание, отвечающее мере $\boldsymbol{P}_{x}$, и где положено $f\left(x_{\tau}\right)=0$ при $\tau \geqslant \zeta$. Другое часто используемое свойство Э. ф. состоит в том, что $P_{x}$ - почти наверное случайная функция $f\left(x_{t}\right)$ непрерывна справа на интервале $[0, \zeta)$ (cM. [3]).



Э. ф. $f(x)<\infty$ наз. г а р м о ни ч е с к о й, если, как бы ни задавался компакт $K \subset E$, выполняется $f(x) \equiv$ $\equiv M_{x} f\left(x_{\tau}\right)$, где $\tau-$ момент первого выхода $X$ из $K$. Потенци а лом наз. любую Ә. ф. $f(x)<\infty$, для $\mathbf{\kappa}-\mathbf{p o u ̈}$



\[

\lim _{n \rightarrow \infty} M_{x} f\left(x_{\tau_{n}}\right)=0

\]



при любом выборе марковских моментов $\tau_{n}, n \geqslant 1$, таких, что $\tau_{n} \rightarrow \zeta$ при $n \rightarrow \infty$. Для части винеровского процесса в области $E \subset \mathbb{R}^{n}$ гармонич. Функции и потенциалы, это, соответственно, неотрицательные в $E$ гармонические в классич. смысле функции и потенциалы Грина борелевских мер, сосредоточенных в $E$.



Примером потенциала служит потенциал $M_{x} \gamma_{\xi}$ аддитивного функционала $\gamma_{t} \geqslant 0$ от $X$, если $M_{x} \gamma_{\zeta}<\infty$. Э. ф. $f(x)<\infty$ является потенциалом аддитивного функционала в том и только том случае, когда



\[

\lim _{n \rightarrow \infty} M_{x} f\left(x_{\tau_{n}}\right) \equiv 0

\]



где $\tau_{n}$ - момент 1-го попадания в множество $\{x: f(x) \geqslant$ $\geqslant n\}$.





\end{document}